% --------------------------------------------------------------------
% 결론 및 질의 섹션입니다.
% Conclusion and Q&A section.
%
% 이 파일에서는 발표 전체를 요약하고 발표 종류에 맞는 질의 쟁점을 정리합니다.
% This file summarizes the whole presentation and lists Q&A issues for the selected presentation stage.
% --------------------------------------------------------------------

% 마지막 큰 섹션입니다. 발표를 닫고 질의응답으로 넘어갑니다.
% Final major section. It closes the presentation and moves into Q&A.
\section{결론 및 질의}

% 요약 하위 섹션입니다.
% Summary subsection.
\subsection{요약}

% 발표 전체를 한 장으로 다시 연결하는 슬라이드입니다.
% This slide reconnects the whole presentation in one place.
\begin{frame}[t]{\BeamerSummaryFrame}
\small

% 핵심 요약은 연구 문제, 이론, 방법, 결과를 같은 언어로 묶어 줍니다.
% The summary ties research questions, theory, methods, and results together.
\begin{block}{핵심 요약}
\begin{enumerate}
\tightitems

% 아래 문장은 발표자의 실제 연구 내용에 맞게 바꿉니다.
% Replace the following sentences with your actual research content.
\item 연구는 \textbf{[핵심 문제]}에서 출발합니다.
\item 이론적 배경은 \textbf{[핵심 개념/분석틀]}로 정리됩니다.
\item 연구 방법은 \textbf{[대상/자료/분석 방법]}을 통해 연구 문제에 답합니다.
\item \KNUEStageText{예상 결과는 \textbf{[산출물/효과/설계 원리]}로 제시됩니다.}{연구 결과는 \textbf{[핵심 발견/해석/시사점]}으로 정리됩니다.}
\end{enumerate}
\end{block}

% 마무리 문장 예시는 발표 마지막 발화를 준비하기 위한 자리입니다.
% The closing sentence example helps prepare the final spoken message.
\begin{exampleblock}{마무리 문장 예시}
본 연구는 \textbf{[연구 대상]}의 \textbf{[핵심 변화 또는 현상]}를 설명하고, \textbf{[교육적 또는 학문적 기여]}를 제안하고자 합니다.
\end{exampleblock}

\slidenote{
마지막 요약은 처음 제시한 연구 목적과 같은 언어로 닫아 줍니다. 발표 전체가 하나의 논리로 이어졌다는 인상을 주는 것이 중요합니다.
}
\end{frame}

% 질의응답 하위 섹션입니다.
% Q&A subsection.
\subsection{\BeamerQASubsection}

% 심사위원에게 어떤 부분의 피드백이나 질의를 받고 싶은지 명시하는 슬라이드입니다.
% This slide states what kind of feedback or questions you want from the committee.
\begin{frame}[t]{\BeamerQAFrame}
\small

% alertblock으로 질의응답의 핵심 쟁점을 강조합니다.
% Use alertblock to emphasize the key issues for Q&A.
\begin{alertblock}{\BeamerQABlock}
\begin{itemize}
\tightitems

% 실제 발표에서는 지도받고 싶은 쟁점을 구체적인 질문으로 바꿉니다.
% In your actual presentation, turn these into specific questions you want advice on.
\item \KNUEStageText{연구 문제의 범위와 표현이 적절한가?}{연구 결과가 연구 문제에 충분히 답하는가?}
\item \KNUEStageText{연구 방법이 각 연구 문제에 답하기에 충분한가?}{결과 해석이 자료와 분석 방법에 근거하는가?}
\item \KNUEStageText{자료 수집 범위와 분석 기준이 현실적인가?}{논의와 시사점이 결과의 범위를 넘어서지 않는가?}
\item \KNUEStageText{예상 결과와 일정이 학위논문으로 완성 가능한 수준인가?}{최종 논문 수정에서 보완할 핵심 쟁점은 무엇인가?}
\end{itemize}
\end{alertblock}

\vspace{0.2cm}

% 감사 인사와 질의응답 안내 문구입니다.
% Closing thanks and Q&A prompt.
\begin{center}
{\Large 감사합니다.}\\[0.2cm]
{\small 질문과 조언 부탁드립니다.}
\end{center}

\slidenote{
\KNUEStageText
{질의응답에서는 방어적으로 답하기보다, 연구계획을 더 정교하게 만들기 위한 검토 의견으로 받아들이는 태도를 유지합니다.}
{질의응답에서는 결과의 한계와 기여를 균형 있게 설명하고, 최종 논문에 반영할 수정 방향을 정리합니다.}
}
\end{frame}
